\chapter{MCMC with a surrogate model}
\label{ch:surrogate}

\section{The idea}

MCMC needs many evaluations --- tens of thousands even in the small examples of \cref{ch:mcmc}.
When each one is an expensive simulation, that can be out of reach even on a cluster. The
alternative is to train a fast \emph{surrogate} (or emulator) of the log-probability on a modest
number of expensive evaluations, and to run the MCMC on the surrogate instead. Once a fast surrogate
exists, the MCMC no longer needs a cluster, or this package: it can be run with any MCMC library,
such as \texttt{emcee} or \texttt{pymc}.

The difficulty is where to put the training data. Sampling the parameter space at random, over a
domain wide enough to be safe, spreads the evaluations thinly. Standard checks of the surrogate,
such as cross-validation on held-out points, then measure its accuracy \emph{on average over that
domain} --- while what matters is its accuracy on the posterior, a region not known in advance and
typically a small fraction of the domain.

This is where a short expensive MCMC is useful even when a full one is unaffordable: its walkers
generate training data exactly where the posterior is. A surrogate-based calibration therefore
proceeds in five steps:

\begin{enumerate}[leftmargin=*, label=\textbf{\arabic*.}]
  \item Start a parallel MCMC with the expensive function.
  \item Run it long enough to generate relevant training data.
  \item Train a fast surrogate on the evaluations, and test its accuracy.
  \item Run the MCMC on the surrogate; no cluster is needed any more.
  \item Validate the resulting posterior against the expensive function, with importance sampling.
\end{enumerate}

\section{Validating with importance weights}
\label{sec:importance}

The surrogate posterior $\psur$ can be checked directly against the expensive one, $\pexp$. Draw
points $\theta_i$ from the surrogate-based MCMC samples, evaluate the expensive function there, and
form the ratios
\begin{equation}
  w_i = \frac{\pexp(\theta_i)}{\psur(\theta_i)} = \exp\!\big(\log\pexp(\theta_i) - \log\psur(\theta_i)\big),
  \label{eq:importance-weights}
\end{equation}
the \textbf{importance weights}. Where the surrogate is exact every ratio is the same constant;
their spread measures the surrogate's error, precisely on the points where the posterior has mass.
The weights also correct the surrogate-based estimate: an expectation under the expensive posterior
is the weighted average over the surrogate samples,
\begin{equation}
  \mathbb{E}_{\pexp}[h(\theta)] \;\approx\; \frac{\sum_i w_i\, h(\theta_i)}{\sum_i w_i},
\end{equation}
self-normalised so that the unknown constants of both densities cancel.

\section{Example: a surrogate for the 2D Rosenbrock}
\label{sec:surrogate-example}

The target is again the constrained 2D Rosenbrock of \cref{sec:mcmc-example}, standing in for an
expensive function. A short expensive MCMC runs with $\Nc = 10$ walkers for $\Lc = 2000$ iterations
--- not converged, which is not the point here. A Gaussian-process surrogate of the log-probability
is then trained on 400 distinct points from that run, plus 200 random points inside the constraint
circle, which regularize the surrogate away from the chains.

\begin{figure}[H]
  \centering
  \includegraphics[width=\linewidth]{example_mcmc_surrogate_2d_visualization.png}
  \caption{The expensive and the surrogate log-probabilities (logarithm of the absolute value),
           with the training points in red. The training data concentrate where they are needed.}
\end{figure}

The MCMC is then rerun, with the same settings, on the surrogate:

\begin{figure}[H]
  \centering
  \includegraphics[width=0.5\linewidth]{example_mcmc_surrogate_parameters_distribution.png}
  \caption{The posterior from the expensive MCMC and from the surrogate-based MCMC, after burn-in.
           They are similar, but not identical.}
\end{figure}

For validation, 1000 points are drawn from the surrogate-based samples and evaluated with the
expensive function:

\begin{figure}[H]
  \centering
  \includegraphics[width=0.42\linewidth]{example_mcmc_surrogate_importance_weights.png}
  \caption{The ratio of expensive to surrogate probabilities, the importance weights of
           \cref{eq:importance-weights}, on 1000 surrogate-posterior points.}
\end{figure}

The ratios are close to 1, to within about 1\%, which validates the surrogate for this problem ---
and the weights could now be used to correct predictions made from the surrogate-based samples.

\begin{keypoint}[From a manual procedure to an automated one]
Every step above required a judgement: how long to run the expensive MCMC, how many regularization
points to add, whether the surrogate is accurate enough, and what to do if it is not. The hybrid
pipeline of \cref{ch:hybrid} iterates the procedure, measures its accuracy on the posterior after
each iteration, recycles the validation evaluations as new training data, and stops when measured
criteria are met.
\end{keypoint}
